

\section{RISC-V Instruction Set Architecture}

The RISC-V Instruction Set Architecture (ISA) is a free, open and modular architecture designed to support both research and industrial applications. Its simplicity, clean encoding, and well-defined extensions mechanisms make it suitable for experimenting with custom accelerators. As this project's focus is on accelerating TM inference through custom processor instructions, the most relevant aspects of the RISC-V ISA are outlined below.

\subsection{RISC-V Overview}

RISC-V follows the classic Reduced Instruction Set Computer (RISC) design philosophy \cite{AsanovicPatterson2014}, emphazising a small set of simple instructions executed efficifently in hardware. The base integer (I) ISA \cite{WatermanLee2011}, RV32I, provides:

\begin{itemize}
    \item 32 general-purpose register
    \item Simple arithmetic and logical operations
    \item Load/store memory access
    \item Branch and jump isntructions
\end{itemize}

A key part of RISC-V is that it uses a modular extension model, where additional instruction groups, or ISA extensions, can be added on top of the base ISA. This is similar to how general-purpose programming languages use libraries to extend the base functionality of the language. Examples of official ISA extensions include:

\begin{itemize}
    \item M: integer multiplication/division
    \item A: atomic operations
    \item F/D: single-precision (F) and double-precision (D) floating point operations
    \item V: vector operations
\end{itemize}

This modularity allows designers to tailor the ISA to the specific workload of their current project, while preserving compatibility with the standard toolchain.

\subsection{Custom Extensions and Instructions}
One of RISC-V’s defining features is its built-in support for custom instructions and custom extensions (sets of instructions). For custom instructions, the ISA reserves opcode space specifically for implementation-defined operations, allowing designers to create application-specific instructions without conflicting with official extensions.

Custom instructions are useful when designing accelerators, like in ML, cryptography or digital signal processing. Some examples of usecases are:

\begin{itemize}
    \item Accelerating computationally heavy loops
    \item Providing efficient access to specialized HW units
    \item Reducing instruction count by combining operations
    \item Supporting algorithm-specific functionality
\end{itemize}

\subsection{Base Instruction Formats}
\label{sec:instform}
Most RISC-V instructions use fixed 32-bit encodings with defined fields for opcode, register operands, immediate values, and function identifiers. There are four core intruction formats in the base ISA \cite{WatermanLee2011}, as shown in Figure \ref{fig:insttypes}.


\begin{figure}[htbp]
\centering
\textbf{R-Type}
\vspace{0.5em}

% ----- R-Type -----
\begin{bytefield}[bitwidth=1.05em]{32}
  \bitheader{31,25,24,20,19,15,14,12,11,7,6,0} \\
  \bitbox{7}{opcode} &
  \bitbox{5}{rd} &
  \bitbox{3}{funct3} &
  \bitbox{5}{rs1} &
  \bitbox{5}{rs2} &
  \bitbox{7}{funct7} 
\end{bytefield}

\vspace{1em}
\textbf{I-Type}
\vspace{0.5em}

% ----- I-Type -----
\begin{bytefield}[bitwidth=1.05em]{32}
  \bitheader{31,20,19,15,14,12,11,7,6,0} \\
  \bitbox{7}{opcode} &
  \bitbox{5}{rd} &
  \bitbox{3}{funct3} &
  \bitbox{5}{rs1} &
  \bitbox{12}{imm[11:0]}  
\end{bytefield}

\vspace{1em}
\textbf{S-Type}
\vspace{0.5em}

\begin{bytefield}[bitwidth=1.05em]{32}
  \bitheader{31,25,24,20,19,15,14,12,11,7,6,0} \\
  \bitbox{7}{opcode} &
  \bitbox{5}{imm[4:0]} &
  \bitbox{3}{funct3} &
  \bitbox{5}{rs1} &
  \bitbox{5}{rs2} &
  \bitbox{7}{imm[11:5]} 
\end{bytefield}

\vspace{1em}
\textbf{U-Type}
\vspace{0.5em}

\begin{bytefield}[bitwidth=1.05em]{32}
  \bitheader{31,12,11,7,6,0} \\
  \bitbox{7}{opcode} &
  \bitbox{5}{rd} &
  \bitbox{20}{imm[19:0]}  
\end{bytefield}

\caption{RISC-V instruction formats (examples).}
\label{fig:insttypes}
\end{figure}


While all types are used in various extensions, we will focus on the R-type, as it offers two source registers, which is useful when implementing custom instructions. From Figure \ref{fig:insttypes} we see that R-type takes in:

\begin{itemize}
    \item \textit{funct7} and \textit{funct3}, which are immediates used to specify functionality
    \item \textit{rs2} and \textit{rs1} are source registers, specifying which register to take values from
    \item \textit{rd} is the destination register, in which the output of the functionality will lie
    \item \textit{opcode} which identifies what instruction group the instruction belongs to (BRANCH, LOAD, STORE, etc.)of length
\end{itemize}
